\section*{Chapitre \ref{ch:b3:plant-molecular-physiology} --- Physiologie végétale moléculaire et réponses au stress}

\begin{solution}{exo:b3:plant-molecular-physiology:1}
\omterm{def:b3:plant-molecular-physiology:photoreceptors}{Phytochromes}, rouge \qty{660}{nm} et rouge lointain \qty{730}{nm}~:
dé-étiolation, \omterm{prop:b3:plant-molecular-physiology:photoequilibrium}{évitement de l'ombre}, germination. \omterm{def:b3:plant-molecular-physiology:photoreceptors}{Cryptochromes}, bleu
\qty{450}{nm}~: dé-étiolation, horloge, floraison. \omterm{def:b3:plant-molecular-physiology:photoreceptors}{Phototropines}, bleu~:
courbure vers la lumière, ouverture des stomates, déplacement des
chloroplastes. La chlorophylle mesure la lumière comme énergie et ne
fait rien de l'information~; les photorécepteurs lisent la couleur, la
direction et la durée à des flux un million de fois plus faibles et
changent l'expression des gènes --- la différence entre manger la
lumière et la lire.
\end{solution}

\begin{solution}{exo:b3:plant-molecular-physiology:2}
Chaque \omterm{def:b3:endocrinology:hormone}{hormone} favorise l'ubiquitination et la destruction d'un
répresseur~: l'\omterm{def:b3:plant-molecular-physiology:hormones}{auxine} colle \omterm{def:b3:plant-molecular-physiology:hormones}{TIR1} aux protéines Aux/IAA, la \omterm{def:b3:plant-molecular-physiology:hormones}{gibbérelline}
liée à GID1 livre les protéines \omterm{def:b3:plant-molecular-physiology:hormones}{DELLA} à une ligase à boîte F, le
jasmonate colle COI1 aux protéines JAZ. Les répresseurs sont déjà posés
sur leurs gènes cibles avec les activateurs prêts à côté, de sorte
qu'aucune protéine nouvelle n'est nécessaire~: détruire le répresseur,
ce que le protéasome fait en quelques minutes, allume les gènes
aussitôt.
\end{solution}

\begin{solution}{exo:b3:plant-molecular-physiology:3}
L'ABA lie PYR/PYL~; le complexe inhibe la phosphatase PP2C~; la kinase
SnRK2 (OST1), qui n'est plus déphosphorylée, devient active et
phosphoryle le canal anionique SLAC1~; les anions sortent, la membrane
se dépolarise, les canaux potassiques sortants s'ouvrent, le K$^{+}$
puis l'eau s'en vont, la turgescence tombe et les \omterm{prop:b3:plant-molecular-physiology:stomata}{cellules de garde}
s'affaissent l'une contre l'autre. Tout maillon peut être rompu~: pas de
récepteurs, une phosphatase qu'on ne peut inhiber (le classique mutant
\emph{abi1}), pas d'OST1 ou pas de SLAC1 --- chacun donne une plante
flétrie dont les stomates restent ouverts dans la sécheresse.
\end{solution}

\begin{solution}{exo:b3:plant-molecular-physiology:4}
La PTI reconnaît des molécules communes à des classes entières de
microbes (flagelline, chitine) par des kinases réceptrices à la surface
de la cellule et se termine par une paroi renforcée, des stomates
fermés, l'expression de gènes de défense et une croissance ralentie,
d'ordinaire sans mort cellulaire. L'ETI reconnaît une protéine
effectrice précise ou le dégât qu'elle fait par un \omterm{def:b3:plant-molecular-physiology:immunity}{récepteur NLR}
intracellulaire et se termine par la mort hypersensible des cellules
infectées et une alarme systémique.
\end{solution}

\begin{solution}{exo:b3:plant-molecular-physiology:5}
Pleine lumière $0.87\times 1.15/1.75 = 0.57$~; une feuille $0.87\times 0.4/1.0 =
0.35$~; couvert dense $0.87\times 0.1/0.7 = 0.12$~; lampe à
incandescence $0.87\times 0.7/1.3 = 0.47$. Une lampe à filament est
riche en rouge lointain, de sorte que $\varphi$ est au-dessous de celui
du jour et que le plant lit une ombre partielle~: il s'allonge. Les
lampes blanc froid et les LED sans rouge lointain font l'inverse.
\end{solution}

\begin{solution}{exo:b3:plant-molecular-physiology:6}
Paroi $5.0$, cytosol $7.5$~: $(1 + 10^{2.75})/(1 + 10^{0.25}) = 563/2.78
= 203$. Paroi $4.5$~: $563/(1 + 10^{-0.25}) = 563/1.56 = 360$ ---
acidifier la paroi élève la fraction neutre au dehors et la capture.
Cytosol $6.5$ avec paroi $5.5$~: $(1 + 10^{1.75})/6.62 = 57/6.62 = 8.6$~:
le piégeage tombe d'un facteur cinq, l'\omterm{def:b3:plant-molecular-physiology:hormones}{auxine} fuit par toutes les faces,
et les gradients polaires s'aplatissent --- une raison pour laquelle les
racines anoxiques perdent leur orientation.
\end{solution}

\begin{solution}{exo:b3:plant-molecular-physiology:7}
\qty{1}{cm/h} $= \qty{10000}{\micro m/h}$~: $200$ cellules par heure,
\qty{18}{s} dans chacune. Diffusion sur \qty{50}{\micro m}~: $t =
(5\times 10^{-5})^{2}/(2\times 5\times 10^{-10}) = \qty{2.5}{s}$.
L'intérieur de la cellule est traversé en quelques secondes~; l'essentiel
des \qty{18}{s} se passe à attendre d'être sorti par les PIN de la
membrane basale, qui est l'étape limitante.
\end{solution}

\begin{solution}{exo:b3:plant-molecular-physiology:8}
$E/A = 1.6\,\Delta w/\Delta c = 1.6\times\num{16000}/150 = 171$ molécules
d'eau par CO$_{2}$. En C$_{4}$, $\Delta c = 250$~: $102$. Jour sec,
$\Delta w = 24$~: $256$.
\end{solution}

\begin{solution}{exo:b3:plant-molecular-physiology:9}
Le gel déshydrate les cellules en tirant l'eau vers la glace
extracellulaire, de sorte que le froid et la sécheresse sont tous deux
des stress de déshydratation. Ils partagent des seconds messagers (ABA,
pointes de calcium, espèces réactives de l'oxygène), des facteurs de
transcription (la famille CBF/DREB est induite par les deux) et des
produits~: déshydrines et protéines LEA qui protègent membranes et
protéines de la perte d'eau, solutés compatibles (proline, sucres) qui
retiennent l'eau, et antioxydants. Une plante qui les a faits pour un
stress les a pour l'autre.
\end{solution}

\begin{solution}{exo:b3:plant-molecular-physiology:10}
À \qty{1}{\%} du plein soleil, $k_{1} + k_{2} = \qty{0.005}{s^{-1}}$~:
95~\% de l'équilibre après $3/k = \qty{600}{s}$, dix minutes, contre six
secondes à midi. Huit heures font quatre demi-vies~: $1/16 \approx
\qty{6}{\%}$ de la Pfr subsiste. Un bref éclair fixe $\varphi$ mais la
Pfr décroît ensuite, alors qu'une journée la tient haute des heures
durant~; les réponses qui exigent que la Pfr agisse des heures
(germination, dé-étiolation) intègrent l'exposition, de sorte qu'une
graine retournée un instant par une charrue ne répond pas comme une
graine posée à la surface.
\end{solution}

\begin{solution}{exo:b3:plant-molecular-physiology:11}
Dix sites par jour à $100$ cellules chacun~: \num{1000} cellules,
$10^{-4}$ de la feuille~; sur une saison une fraction de pour cent,
contre la perte de la feuille si une infection se répandait sans
obstacle. Les cellules mortes ne coûtent presque rien et emportent avec
elles la nourriture du pathogène --- pour un biotrophe, qui a besoin de
cellules vivantes. Un nécrotrophe se nourrit de tissu mort~: la \omterm{def:b3:plant-molecular-physiology:immunity}{réaction
hypersensible} lui offre un repas et un pied à terre, et certains
(\emph{Botrytis}, \emph{Sclerotinia}) sécrètent des toxines qui la
provoquent exprès~; contre eux la plante s'appuie sur les défenses
médiées par le jasmonate et non sur la mort cellulaire.
\end{solution}

\begin{solution}{exo:b3:plant-molecular-physiology:12}
Production au taux $p$ de la douzième à la seizième heure. Jour long,
lumière tout du long, $k = \ln 2/4 = \qty{0.17}{h^{-1}}$~: le niveau à
16 vaut $(p/k)(1 -
\mathrm{e}^{-4k}) = (p/0.17)(0.5) = 2.9p$. Jour court, obscurité dès la
dixième heure, $k = \qty{1.39}{h^{-1}}$~: $(p/1.39)(1 - \mathrm{e}^{-5.5})
\approx 0.72p$. Quatre fois plus de CONSTANS dans le jour long. Un seuil
entre les deux --- disons $2p$, nécessaire pour activer \emph{FT} ---
n'est franchi que lorsque la fenêtre de production de l'horloge coïncide
avec la lumière~: la coïncidence d'un rythme interne avec le jour
externe convertit une durée graduée en une décision tout ou rien de
fleurir.
\end{solution}

\begin{solution}{pb:b3:plant-molecular-physiology:1}
\textbf{1.} Pleine lumière $0.87\times 1.15/1.75 = 0.57$~; couvert $0.87\times
0.2/0.8 = 0.22$.
\textbf{2.} $r = 2(1 - \varphi)$~: \qty{0.86}{mm/h} en pleine lumière,
\qty{1.57}{mm/h} à l'ombre~; sur \qty{48}{h}~: $41$ contre \qty{75}{mm},
soit \qty{34}{mm} de plus.
\textbf{3.} Le rouge tombe à $0.1\times 1.15 = 0.115$ du rouge lointain,
qui tombe à $0.8$~: $\zeta = 0.144$, $\varphi = 0.87\times 0.144/0.744 =
0.17$. Oui~: une seule feuille fait tomber $\varphi$ de $0.57$ à $0.17$,
en plein dans la plage de l'\omterm{prop:b3:plant-molecular-physiology:photoequilibrium}{évitement de l'ombre}.
\textbf{4.} Les deux vitesses de photoconversion varient avec le flux,
donc leur rapport non. Sous les nuages le spectre est presque inchangé
et $\varphi$ reste voisin de $0.57$~: un plant à découvert par temps
gris ne s'allonge pas --- il lit les voisines, non l'intensité.
\textbf{5.} \qty{95}{\%} après $3/(k_{1} + k_{2})$~: \qty{6}{s} en plein
soleil, \qty{600}{s} à \qty{1}{\%}.
\textbf{6.} \qty{8}{h} $=$ quatre demi-vies, $1/16 = \qty{6.3}{\%}$~;
\qty{14}{h} $=$ sept, $1/128 = \qty{0.8}{\%}$. La Pfr résiduelle à
l'aube pourrait rapporter la longueur de la nuit, mais le retour en
arrière est une réaction thermique qui court plus vite les nuits
chaudes, et la valeur de départ dépend de la lumière du jour~: les
plantes mesurent plutôt la longueur de la nuit avec l'horloge.
\textbf{7.} Paroi~: $1/(1 + 10^{0.75}) = 0.15$~; cytosol~: $1/(1 +
10^{2.45}) = 0.0035$.
\textbf{8.} $(1 + 282)/(1 + 5.6) = 283/6.6 = 43$. Cytosol
$43\times 0.1 = \qty{4.3}{\micro M}$.
\textbf{9.} $\qty{1}{cm/h}/\qty{100}{\micro m} = 100$ cellules par heure,
\qty{36}{s} chacune.
\textbf{10.} \qty{5}{h}. La courbure commence en moins d'une heure parce
que la redistribution qui compte est latérale, en travers d'un
millimètre de pointe, et la réponse est locale~; le courant à longue
distance fixe l'approvisionnement, non le tempo.
\textbf{11.} Flanc inférieur~: $60/50 = 1.2$, \omterm{def:b3:plant-molecular-physiology:hormones}{auxine} \qty{20}{\%}
au-dessus de la norme, élongation $0.8$~; flanc supérieur~: $0.8$ de la
norme, élongation $1.2$. Le flanc supérieur pousse plus que l'inférieur~:
la racine s'incurve vers le bas.
\textbf{12.} Les cellules de tige sont au-dessous de leur optimum
d'\omterm{def:b3:plant-molecular-physiology:hormones}{auxine}, de sorte que l'\omterm{def:b3:plant-molecular-physiology:hormones}{auxine} en plus du flanc inférieur favorise leur
croissance~: le flanc inférieur pousse plus que le supérieur et la tige
s'incurve vers le haut.
\textbf{13.} $E = 0.3\times 16 = \qty{4.8}{mmol\,m^{-2}\,s^{-1}}$.
\textbf{14.} $A = (0.3/1.6)\times 150 = \qty{28}{\micro mol\,m^{-2}\,s^{-1}}$~;
$E/A = 4800/28 = 171$.
\textbf{15.} Surface \qty{0.002}{m^2}, \num{43200}~s~: eau $4.8\times
10^{-3}\times 0.002\times\num{43200} = \qty{0.41}{mol} = \qty{7.5}{g}$~;
CO$_{2}$ $28\times 10^{-6}\times 0.002\times\num{43200} =
\qty{2.4e-3}{mol}$, soit $2.4\times 10^{-3}/6\times 180 =
\qty{0.073}{g}$ de glucose --- cent grammes d'eau pour un gramme de
sucre.
\textbf{16.} $E = 0.05\times 16 = \qty{0.8}{mmol\,m^{-2}\,s^{-1}}$,
$A = \qty{4.7}{\micro mol\,m^{-2}\,s^{-1}}$, rapport toujours $171$. La
plante a épargné six septièmes de son eau et renoncé à six septièmes de
son carbone~: fermer la soupape change le débit, non le prix.
\textbf{17.} C$_{4}$~: $\Delta c = 250$, $E/A = 4800/46.9 = 102$. CAM la
nuit~: $E = 0.3\times 5 = \qty{1.5}{mmol}$, $A = 28$~: $E/A = 53$.
\textbf{18.} $E/A = 1.6\,\Delta w/\Delta c$~: la conductance s'élimine
parce que les deux gaz passent par le même pore. La plante peut changer
les gradients --- concentrer le CO$_{2}$ au dedans (C$_{4}$), ouvrir
quand l'air est humide (CAM, aube), refroidir la feuille, épaissir la
couche limite avec des poils --- mais non la physique de deux gaz
partageant un seul trou.
\textbf{19.} \num{1000} cellules par jour, $10^{-4}$ de la feuille.
\textbf{20.} \num{60000} cellules en 60 jours, \qty{0.6}{\%}. Une
échappée par jour détruisant \qty{5}{\%} consommerait la feuille en
vingt jours~: la \omterm{def:b3:plant-molecular-physiology:immunity}{réaction hypersensible} coûte le centième de ce que
coûte une seule infection échappée.
\textbf{21.} $0.2\times 5 = \qty{1}{\%}$ de la photosynthèse. Laissée
allumée en permanence elle coûterait \qty{5}{\%} et la plante serait
dépassée par des voisines qui mettent cela dans des feuilles~; une
défense induite ne paie que lorsque la menace est là.
\textbf{22.} Il y gagne l'invisibilité à ce NLR et infecte la variété
résistante~; il y perd ce que faisait l'effecteur (étouffer la PTI), de
sorte qu'il est plus faible sur les plantes sans le NLR. La population
végétale en vient à favoriser des NLR contre les effecteurs restants, et
le vieux gène de résistance, devenu inutile, décline~: un cycle dépendant
de la fréquence des gènes des deux camps.
\textbf{23.} L'imitateur (la coronatine) active la voie du jasmonate,
qui antagonise la voie du salicylate~; le salicylate est la défense
contre les biotrophes tels que la bactérie elle-même, et l'imitateur
rouvre aussi les stomates que la plante avait fermés. Le pathogène
retourne un bras du système immunitaire contre l'autre.
\textbf{24.} Un gène de résistance unique sur des millions de plantes
identiques est une sélection uniforme~: tout mutant qui perd ou modifie
l'effecteur se répand sur la culture entière en quelques saisons. Les
sélectionneurs empilent plusieurs gènes de résistance pour qu'il faille
perdre plusieurs effecteurs à la fois, mélangent les variétés, font
tourner les gènes d'une année à l'autre, et favorisent une large
résistance déclenchée par les motifs qu'aucune mutation unique n'esquive.
\textbf{25.} $\varphi \approx 0.22$ sous le couvert~; \omterm{def:b3:plant-molecular-physiology:hormones}{auxine} dedans :
dehors $\approx 43$~; environ $170$ molécules d'eau perdues par CO$_{2}$
fixé.
\end{solution}
